\documentclass{UVKABoA4}

\usepackage{graphicx} % Standardpaket zur Grafikeinbindung
\usepackage{amsmath,amssymb} % Erweiterung des Mathematik-Modus
\usepackage[colorinlistoftodos, german]{todonotes} % Option 'disable' entfernt alle ToDos
\usepackage[absolute,overlay]{textpos}
\usepackage{vmargin}          % Adjust margins in a simple way
\usepackage{tikz}

\usepackage{listings}
\lstset{numbers=left,					% where to put the line-numbers
	language=C++,								% choose the language of the code
	basicstyle=\ttfamily\small\mdseries,	% the size of the fonts that are
	% used for the code; \ttfamily\scriptsize\mdseries, \small, \footnotesize
	numberstyle=\tiny,					% the size of the fonts that are used for the line-numbers; \tiny, \footnotesize
	stepnumber=1,								% the step between two line-numbers. If it's 1 each line will be numbered
	numbersep=5pt,							% how far the line-numbers are from the code
	backgroundcolor=\color{hellgrau},	% choose the background color. You must add \usepackage{color}
	showspaces=false,						% show spaces adding particular underscores
	tabsize=2,									% sets default tabsize to 2 spaces
	captionpos=b,								% sets the caption-position to bottom
	breaklines=true,						% sets automatic line breaking
	title=\lstname,							% show the filename of files included with \lstinputlisting; also try caption instead of title
	keywordstyle=\color{blue},
	xleftmargin=30pt,
	xrightmargin=10pt,
	aboveskip=\baselineskip
}%
% \lstloadlanguages{C++}%


% Toggle the following two lines to switch between english and german layout
% \usepackage[english]{babel}	% Englische Rechtschreibung und Silbentrennung.
\usepackage[ngerman]{babel}		% Neue deutsche Rechtschreibung und Silbentrennung.

\usepackage[raiselinks=true,
            bookmarks=true,
            bookmarksopenlevel=1,
            bookmarksopen=true,
            bookmarksnumbered=true,
            hyperindex=true,
            plainpages=false,
            pdfpagelabels=true,
            pdfborder={0 0 0.5}]{hyperref}

% Glossar und Akronyme
% Übersetzung:
%		1. pdflatex
%		2. makeglossaries
%		3. pdflatex
\usepackage[ngerman]{translator}
\usepackage[toc, footnote]{glossaries}		% [acronym] -option legt ein Extra-Verzeichnis für Abkürzungen an
\makeglossaries


\renewcommand{\labelitemi}{$\triangleright$} 	% Aufzählungszeichen verändern
\newcommand{\authorsection}[1]{\begin{flushright}\textit{Autor: #1}\end{flushright}}
\newcommand{\todoprivate}[1]{\todo[color=green]{#1}}

\makeatletter
\def\tagform@#1{\maketag@@@{\ignorespaces [#1]\unskip\@@italiccorr}}	% Anpassung der Formelnummerierung
\makeatother

\makeindex

\newcommand{\editorabel}{Abel~Gawel}
\newcommand{\editoranne}{Anne~Günter}
\newcommand{\editorhamza}{Hamza~Hajj~Hassan}
\newcommand{\editormarton}{M\'{a}rton~Kelemen}
\newcommand{\editoroier}{Oier~Mees}
\newcommand{\editortobias}{Tobias~Roth}
\newcommand{\editordirk}{Dirk~Toewe}
\newcommand{\editorandreas}{Andreas~Weber}
\newcommand{\editorjulian}{Julian~Wormer}

\newcommand{\editordummy}{Dummy~Dummy}

\newcommand{\headtitle}{Projektpraktikum mit mobilen Robotern}
\newcommand{\floatingtitle}{Projektpraktikum mit mobilen Robotern}
\newcommand{\thesistype}{Ausarbeitung}

\newcommand{\advisorone}	{Dipl.--Inform.~Andreas~Hermann}
\newcommand{\advisortwo}	{Dipl.--Inform.~Jan~Oberländer~(M.\,Sc.)}
\newcommand{\advisorthree}	{Dipl.--Inform.~Zhixhing~Xue}

\newcommand{\timestart}{Oktober 2011}
\newcommand{\releaseyear}{2012}
\newcommand{\timeend}{April \releaseyear}
\newcommand{\releasemonth}{im April \releaseyear}


% Names of the organizations
\newcommand{\department}{\iflanguage{english}{Department of Computer Science}
                                             {Fakultät für Informatik}}
\newcommand{\institute}{\iflanguage{english}{Institute for Anthropomatics}
                                            {Institut für Anthropomatik}}
\newcommand{\fzidepartment}{Abteilung Technisch Kognitive Assistenzsysteme}
\newcommand{\fziname}{\iflanguage{english}{FZI Research Center for Information Technology}
                                          {FZI Forschungszentrum Informatik}}

% A graphic best describing your thesis
\newcommand{\titlefig}{\includegraphics[width=0.4\textwidth]{graphics/FZI-Logo}}

